\section{Căn cứ pháp lý và quy chuẩn dữ liệu y tế Việt Nam}
\label{sec:legal_regulations}

Mọi quyết định lâm sàng, quy tắc kiểm tra tương tác, điều kiện thanh toán bảo hiểm và phân hệ xuất bản đơn thuốc điện tử trong đề tài được xây dựng tuân thủ nghiêm ngặt các văn bản quy phạm pháp luật chính thức của Quốc hội, Bộ Y tế và Bảo hiểm Xã hội Việt Nam:

\begin{itemize}[leftmargin=1.5em]
    \item \textbf{Luật Khám bệnh, chữa bệnh số 15/2023/QH15 \cite{luat_kbcb_2023}:} Khung pháp lý tối cao điều chỉnh hoạt động y tế tại Việt Nam (có hiệu lực từ ngày 01/01/2024), quy định bắt buộc về ứng dụng công nghệ thông tin, bảo vệ an toàn thông tin hồ sơ bệnh án điện tử, trách nhiệm bảo đảm an toàn người bệnh và kiểm soát kê đơn hợp lý.
    \item \textbf{Thông tư số 26/2025/TT-BYT (Ngày 20/08/2025) \cite{tt26_2025_byt}:} Quy định chi tiết về hoạt động kê đơn thuốc và đơn thuốc điện tử tại các cơ sở khám bệnh, chữa bệnh. Thông tư chuẩn hóa 3 mẫu đơn thuốc quốc gia: Đơn thuốc thông thường (Mẫu 1), Đơn thuốc "N" cho thuốc gây nghiện và giảm đau trung ương (Mẫu 2), Đơn thuốc "H" cho thuốc hướng tâm thần và tiền chất (Mẫu 3); đồng thời bắt buộc ghi rõ số định danh định mức ngày dùng tối đa (không quá 30 ngày cho bệnh mạn tính, 7--10 ngày cho thuốc gây nghiện) và kết nối Cơ sở Dữ liệu Đơn thuốc Quốc gia.
    \item \textbf{Dược thư Quốc gia Việt Nam Lần thứ III (Ban hành năm 2022) \cite{dtqg_2022}:} Văn bản chuẩn mực cao nhất về tiêu chuẩn sử dụng thuốc tại Việt Nam, quy định liều dùng chính thức, chỉ định, chống chỉ định, thận trọng khi dùng thuốc cho phụ nữ có thai/cho con bú, và bảng chỉnh liều theo thanh thải Creatinine cho 1.400 chuyên luận hoạt chất.
    \item \textbf{Quyết định số 5948/QĐ-BYT (Ngày 30/12/2021) \cite{qd5948_byt}:} Ban hành "Danh mục tương tác thuốc chống chỉ định trong thực hành lâm sàng tại các cơ sở khám bệnh, chữa bệnh", quy định các cặp hoạt chất tuyệt đối không được kê đồng thời do nguy cơ gây tử vong hoặc tai biến nghiêm trọng (ví dụ: Clopidogrel + Omeprazole/Esomeprazole; Simvastatin + Clarithromycin).
    \item \textbf{Thông tư số 20/2022/TT-BYT (Ngày 31/12/2022) \cite{tt20_2022_byt}:} Ban hành Danh mục và tỷ lệ, điều kiện thanh toán đối với thuốc hóa dược, sinh phẩm thuộc phạm vi được hưởng của người tham gia BHYT, phân định rõ phạm vi thuốc được kê tại bệnh viện Hạng Đặc biệt, Hạng I, Hạng II, Hạng III và Hạng IV nhằm ngăn ngừa xuất toán bảo hiểm.
    \item \textbf{Quyết định số 1666/QĐ-BHXH (Ngày 03/12/2020) \cite{qd1666_bhxh}:} Ban hành mẫu thẻ bảo hiểm y tế với cấu trúc mã thẻ 15 ký tự, quy định mã mức hưởng quyền lợi bảo hiểm (ký tự thứ 4 từ mức 1 đến 5 tương ứng chi trả 80\% đến 100\%) và mã cơ sở khám chữa bệnh ban đầu.
\end{itemize}
